
This study set out to test whether in-distribution trustworthiness benchmark rankings predict out-of-distribution rankings, and whether the answer differs by dimension --- specifically, whether adversarial benchmark design disrupts rank stability for robustness but not for fairness.

The mechanism was falsified. Both fairness and adversarial robustness rankings are highly stable after controlling for general capability: near-perfect for fairness (partial Spearman $\rho$ = 0.962, p < 0.0001, N = 16), substantial for robustness ($\rho_{ANLI}$ = 0.684, p = 0.007; $\rho_{OOD}$ = 0.868, p = 0.0001; N = 13), with zero rank reversals across all adversarial pairs. The predicted disruption does not occur at scale across diverse model families.

A directional advantage for fairness is observed ($\Delta\rho$ = 0.192, Fisher z p = 0.024, N = 13), but the preregistered magnitude criterion ($\Delta\rho \geq$ 0.200) is not met, and the mechanism underlying the directional advantage is unknown.

\textbf{Contributions:}

\begin{enumerate}
\item \textbf{Empirical baseline.} First computation of partial Spearman $\rho$ (MMLU-controlled) between in-distribution and OOD trustworthiness benchmark rankings across 16 LLMs, establishing that both fairness and adversarial robustness rankings are highly stable ($\rho \geq$ 0.68).
\end{enumerate}

\begin{enumerate}
\item \textbf{Methodological demonstration.} MMLU-controlled partial Spearman $\rho$ is a tractable, data-efficient operationalization of trustworthiness benchmark predictive validity using published scores only.
\end{enumerate}

\begin{enumerate}
\item \textbf{Theoretical revision.} The adversarial disruption hypothesis is falsified. Both dimensions reflect stable latent model properties, with fairness showing a directional advantage whose mechanism remains an open question.
\end{enumerate}

\textbf{Future directions.} (1) BBQ item-level Jaccard analysis to distinguish latent model mechanism from benchmark structural artifact --- the highest-priority open question given the alternative hypotheses in Section 6.2. (2) Richer capability covariates (BIG-Bench, GSM8K) to test sensitivity of the $\Delta\rho$ finding beyond MMLU. (3) Instruction-tuning effects: whether RLHF alignment differentially improves fairness generalization relative to robustness generalization remains the most practically actionable question.

Knowing that a model's trustworthiness ranking is stable across distribution shifts is a prerequisite for principled benchmark-based model selection. This study establishes that the prerequisite is empirically met for the model population and benchmarks examined, while identifying the mechanism of any fairness-robustness differential as an open research problem.

